\chapter{河牌 / River}

\section{本章导入}

本章讨论河牌圈。河牌是最终决策街，所有听牌要么完成，要么落空，双方都不再拥有未来街次的补偿空间。因此，河牌决策必须更加明确：是价值下注、诈唬下注、抓诈唬，还是放弃摊牌。

新手在河牌最容易被结果情绪影响：害怕被反超而不敢价值下注，或因为不甘心而用错误价格跟注。本章将帮助读者建立更稳定的河牌判断标准。

\section{核心概念}

本章将介绍摊牌价值、价值下注、薄价值下注、诈唬下注、抓诈唬、阻断牌和面对大注的决策。读者需要理解，河牌行动的依据来自范围、赔率和对手倾向，而不是单纯的害怕或好奇。

\section{新手常见误区}

常见误区包括错过价值下注、用中等牌力下注后只能被更好牌跟注、面对大注过度好奇跟注，以及在没有弃牌率的对手面前强行诈唬。

\section{实战思考框架}

河牌可以先判断自己的牌是否能被更差牌跟注，再判断对手是否会弃掉更好牌。面对下注时，则应结合底池赔率、对手价值范围和潜在诈唬组合做决定。

\section{牌例拆解}

\subsection{牌例一：待补充}

本牌例将用于分析河牌薄价值下注的边界。

\subsection{牌例二：待补充}

本牌例将用于分析面对河牌大注时的抓诈唬条件。

\section{本章总结}

河牌要求读者用最清晰的标准做最终选择。好的河牌决策不是勇敢或保守，而是价格、范围和对手倾向共同支持的结果。

\section{课后训练}

请复盘三手河牌面对下注的牌，计算当时跟注所需胜率，并判断自己的跟注是否有数学依据。
